\documentclass[11pt,a4paper]{article}
\usepackage{common}
\setrunninghead{Microelectronics Laboratory ECC-353 -- CBS Experiment 11}
\hypersetup{pdftitle={ECC-353 CBS Experiment 11 -- Break-Before-Make CMOS Inverter},pdfauthor={}}
\newcommand{\sect}[1]{\section*{#1}}

\begin{document}

\collegeheader

\begin{center}
  {\bfseries MICROELECTRONICS LABORATORY -- ECC-353: CONTENT BEYOND SYLLABUS (CBS)}\\[3pt]
  \textbf{Academic Year:} 2026--2027\\[6pt]
  {\bfseries EXPERIMENT NO.\ 11: Laboratory CBS: Implementation and Simulation of the ECC-305 Selected Circuit}
\end{center}

\studentdetails
\evaluatorbox

\vspace{4pt}
\section*{\centering TITLE OF EXPERIMENT}
\begin{center}\textbf{Implementation and Simulation of a Break-Before-Make CMOS Inverter for Low-Power Delay Generation}\end{center}

\section*{\centering CIRCUIT SELECTED FROM ECC-305}
\textbf{Circuit Name:} Break-Before-Make (BBM) CMOS inverter / delay cell (PtNf full-swing topology, Fig.~9(b) and Fig.~8 of the paper), compared with a conventional 4-inverter CMOS delay cell.

\textbf{Low-Power Technique:} Break-before-make (dual-ramp, time-skewed) switching to eliminate direct-path (short-circuit) current; supply charge/energy per delay is reduced.

\newpage
% ------------------------------------------------------------------
\sect{1. AIM}
To implement and simulate the selected CMOS circuit using the low-power design technique studied during the Microelectronics Theory (ECC-305) research activities and to evaluate its performance.

\sect{2. OBJECTIVES}
\begin{itemize}
  \item To understand the transistor-level architecture.
  \item To implement the selected CMOS circuit.
  \item To verify its functional operation.
  \item To simulate the circuit using the prescribed tool.
  \item To measure power and delay.
  \item To calculate PDP.
  \item To compare the reference and optimized circuit.
  \item To evaluate the effectiveness of the selected/implemented low-power technique(s).
\end{itemize}

\sect{3. RESEARCH REFERENCE}
\textbf{Research Paper Title:} Break-before-Make CMOS Inverter for Power-Efficient Delay Implementation

\textbf{Author(s):} J.~Puhan, D.~Rai\v{c}, T.~Tuma, and \'A.~B\H{u}rmen

\textbf{Journal/Conference:} The Scientific World Journal, vol.~2014, Article ID 349131 (doi:10.1155/2014/349131)

\textbf{Year:} 2014

\sect{4. SOFTWARE / SIMULATION ENVIRONMENT}
\textbf{Simulation Tool:} ngspice-42 (open-source SPICE, batch mode, Ubuntu~24 container). Transient analysis; Python~3 / NumPy used only for post-processing and cross-check. \emph{Not} the paper's HSPICE.

\textbf{Technology Node / Model:} \emph{Substitute} Level-1 (Shichman--Hodges) \SI{0.18}{\micro\meter}-class model built from the parameters in the caption of Fig.~4 of the paper ($V_{TN}=\SI{0.5}{\volt}$, $V_{TP}=\SI{-0.5}{\volt}$, $K_{PN}=\SI{350}{\micro\ampere\per\volt\squared}$, $K_{PP}=\SI{90}{\micro\ampere\per\volt\squared}$, $C_{ox}=\SI{8.6}{\femto\farad\per\micro\meter\squared}$, $L_{\min}=\SI{180}{\nano\meter}$). It is \emph{not} the TSMC \SI{0.18}{\micro\meter} HSPICE model used by the authors (proprietary); $\lambda$, $\gamma$, $\phi$, junction and overlap capacitances are my assumptions (file \texttt{models\_180nm\_L1.lib}).

\textbf{\VDD:} \SI{1.8}{\volt} (same for both circuits)

\textbf{Temperature:} \SI{27}{\celsius} (ngspice \texttt{.temp 27}; Level-1 parameters are not temperature-tuned)

\textbf{Load Capacitance:} $C_L=\SI{2}{\femto\farad}$ at the output of the output unit inverter ($\approx$ one unit-inverter input capacitance); a sweep over \SIrange{0.5}{10}{\femto\farad} was also run. Input: \SIrange{0}{1.8}{\volt} pulse, \SI{50}{\pico\second} edges, \SI{20}{\nano\second} period (\SI{50}{\mega\hertz}), \SI{50}{\percent} duty.

\sect{5. CIRCUIT DIAGRAM}
Transistor-level schematics of the two circuits exactly as simulated (redrawn from paper Figs.~8 and 9; sizes $W/L$ in nm).

\begin{figure}[H]
  \centering
  \includegraphics[width=\textwidth]{fig_circuit_conventional.png}
  \caption{Reference: conventional delay cell (4 static CMOS inverters). Netlist: \texttt{conventional\_inverter.sp}.}
\end{figure}

\begin{figure}[H]
  \centering
  \includegraphics[width=\textwidth]{fig_circuit_bbm_cell.png}
  \caption{Optimised circuit: BBM delay cell, PtNf topology. Netlist: \texttt{bbm\_inverter.sp} (sizes shown are the paper's \SI{750}{\pico\second} sizing).}
\end{figure}

\sect{6. DESIGN / IMPLEMENTATION DETAILS}
\textbf{Circuit reconstruction.} The paper's Fig.~9(b) delay cell is: input unit inverter $\to$ split-output BBM inverter (CMOS-to-BBM converter, Fig.~1(a)) $\to$ split-input merge inverter (BBM-to-CMOS converter, Fig.~1(c)) $\to$ output unit inverter. The bidirectional delay element is the PtNf full-swing implementation of Fig.~8: $P_1$ (source \VDD), node A ($p_o$), $P_t$ (PMOS, gate tied to GND, i.e.\ always conducting in triode), node B ($n_o$), $N_1$ (source GND); $N_f$ is an NMOS connected across A--B whose gate is driven by the next-stage output ($d_o$) to restore the low level of A. Stage~3 has $P_2$ (gate A) and $N_2$ (gate B) with joined drains ($d_o$). These connections were read directly from the figures in the PDF; the polarity of $N_f$'s gate (connected to the merged output) and the two-transistor stage~3 follow the paper's statement that the level-restoration control signals are taken from the next-stage output. Transistor count: conventional 8; BBM 10 (two unit inverters $=4$, plus $P_1,P_t,N_1,N_f=4$, plus $P_2,N_2=2$).

\textbf{Sizing used.} Both cells use the sizing printed in the Fig.~11 caption of the paper for $T_p=\SI{750}{\pico\second}$. The unit inverters are not specified by the paper; I chose $W_p/W_n=\SI{860}{\nano\meter}/\SI{220}{\nano\meter}$, $L=\SI{180}{\nano\meter}$, following the symmetric stage of Fig.~4. Conventional: inv1 $W_p/L_p=220/370$, $W_n/L_n=220/880$; inv2 $W_p/L_p=880/730$, $W_n/L_n=1050/1320$ (nm). BBM: $P_1$ 220/740, $N_1$ 390/180, $P_t$ 220/350, $N_f$ 220/950, $P_2$ 220/180, $N_2$ 220/220 (nm). Source/drain parasitics follow the paper's rule $A_d=A_s=\SI{0.8}{\micro\meter}\times w$, $P_d=P_s=\SI{1.5}{\micro\meter}+w$.

\textbf{Measurement definitions (identical for both circuits).} Delay: \SI{50}{\percent} (\SI{0.9}{\volt}) crossing of the input to \SI{50}{\percent} crossing of the cell output, rising and falling, third edge of each (steady state), $T_p=(t_{pLH}+t_{pHL})/2$. Charge: supply charge of the whole cell (all four stages; the input source is not counted) integrated over 4 complete cycles (\SIrange{21}{101}{\nano\second}) and divided by 4, giving $Q$ per cycle (one rising plus one falling slope, the paper's metric). Energy per cycle $E=Q\,V_{DD}$; average power $P_{avg}=E/T_{cycle}$ at \SI{50}{\mega\hertz}; $\mathrm{PDP}=P_{avg}\times T_p$. Because $E$ per cycle is independent of frequency (leakage is zero in a Level-1 model), $P_{avg}$ and PDP scale with the chosen \SI{50}{\mega\hertz} test frequency; $Q$ (or $E$) per cycle is the frequency-independent comparison. Area is the gate-area proxy $\sum W\cdot L$ over all transistors (not a layout area).

\textbf{Tools and flow.} \texttt{run\_all.py} (baseline, scaling sweeps, $C_L$ and edge-rate sweeps, no-skew ablation); \texttt{optimize.py} (differential-evolution re-sizing at $T_p$ targets \SI{0.5}{\nano\second} and \SI{1.0}{\nano\second}; constraints: $|T_p/\mathrm{target}-1|\le\SI{2}{\percent}$, $|t_{pLH}-t_{pHL}|\le\SI{15}{\percent}$ of $T_p$, output swing $\ge\SI{1.7}{\volt}$ / $\le\SI{0.1}{\volt}$, BBM node levels restored and isolation times $>0$); \texttt{validate\_opt.py} (re-runs the optimised sizes on the full testbench). Checks performed: (i)~convergence errors (none in the final runs); (ii)~no floating nodes (all nodes have DC paths; the DC operating point converged); (iii)~independent Python trapezoidal integration of the raw waveform reproduces the ngspice charge to $<\SI{0.01}{\percent}$ (\SI{67.349}{\femto\coulomb} vs \SI{67.354}{\femto\coulomb} and \SI{40.304}{\femto\coulomb} vs \SI{40.310}{\femto\coulomb}); (iv)~one bug turned up and was fixed: ngspice treats the name \texttt{gnd} as node~0, which silently shorted a current-sense source.

\sect{7. FUNCTIONAL VERIFICATION}
The following are provided in this section:
\begin{itemize}
  \item Input waveform
  \item Output waveform
  \item Truth table, where applicable
  \item Functional verification
\end{itemize}

\textbf{Observation.} Input and output waveforms of both circuits over one steady-state cycle (data from \texttt{conv\_wave.csv} / \texttt{bbm\_wave.csv} produced by ngspice) are in Fig.~\ref{fig:inout}; the BBM internal skewed nodes are in Fig.~\ref{fig:skew}. Both circuits are non-inverting overall (4 inversions).

\begin{figure}[H]
  \centering
  \includegraphics[width=0.78\textwidth]{fig_wave_in_out.png}
  \caption{Functional verification: input (dashed) and output for the conventional and BBM cells (paper's \SI{750}{\pico\second} sizing).}
  \label{fig:inout}
\end{figure}

\begin{figure}[H]
  \centering
  \includegraphics[width=0.62\textwidth]{fig_wave_bbm_skew.png}
  \caption{BBM mechanism: node A (PMOS gate of stage~3) and node B (NMOS gate of stage~3) switch at different times; the stage~3 output $d_o$ follows.}
  \label{fig:skew}
\end{figure}

\begin{table}[H]
  \centering\small
  \begin{tabular}{@{}c c c c@{}}
    \toprule
    in & out (ideal) & out (conventional, simulated steady state) & out (BBM, simulated steady state)\\
    \midrule
    0 & 0 & \SI{-0.014}{\volt} (logic 0) & \SI{-0.009}{\volt} (logic 0)\\
    1 (\SI{1.8}{\volt}) & 1 & \SI{1.827}{\volt} (logic 1) & \SI{1.808}{\volt} (logic 1)\\
    \bottomrule
  \end{tabular}\\[3pt]
  {\footnotesize\emph{The table shows extreme values over the measured window, including small over/undershoot; settled levels are \SI{0}{\volt} and \SI{1.8}{\volt}.}}
\end{table}

\textbf{Observation.} Both circuits reproduce the input with a delay of about \SIrange{0.5}{0.6}{\nano\second} and full-swing outputs ($V_{out,\max}=\SI{1.83}{\volt}$ / \SI{1.81}{\volt}, $V_{out,\min}=\SI{-0.014}{\volt}$ / \SI{-0.009}{\volt} for conventional / BBM). In the BBM cell, on a rising input, node A (PMOS gate) rises \SI{155}{\pico\second} before node B (NMOS gate), and on a falling input B falls \SI{258}{\pico\second} before A (\SI{50}{\percent} crossings); an isolation interval therefore exists in every transition and the merge stage never has both gates in the conducting range at once. Node A overshoots to \SI{2.04}{\volt} and node B undershoots to \SI{-0.10}{\volt} through gate-drain (Miller) coupling and then settle to the full rails ($N_f$ restores the low level of A). Function verified.

\sect{8. PERFORMANCE MEASUREMENT}
\emph{Column meaning:} \textbf{Reference/Existing} is the conventional 4-inverter cell, re-sized (my own differential-evolution optimisation) for $T_p\approx\SI{1.0}{\nano\second}$. \textbf{Student Circuit} is my BBM cell exactly as built from the paper's published \SI{750}{\pico\second} sizing, not re-sized; its $T_p=\SI{563}{\pico\second}$ does \emph{not} match the reference, so it is shown for completeness only and must not be used for an improvement figure. \textbf{Optimized Circuit} is the BBM PtNf re-sized by the same optimiser for $T_p\approx\SI{1.0}{\nano\second}$. All values are from ngspice, full testbench, \SI{27}{\celsius}, $C_L=\SI{2}{\femto\farad}$.

\begin{table}[H]
  \centering\small
  \begin{tabular}{@{}l r r r@{}}
    \toprule
    Parameter & Reference/Existing & Student Circuit & Optimized Circuit\\
    \midrule
    \VDD & \SI{1.8}{\volt} & \SI{1.8}{\volt} & \SI{1.8}{\volt}\\
    Transistor Count & 8 & 10 & 10\\
    Power & \SI{5.84}{\micro\watt} & \SI{3.63}{\micro\watt} & \SI{4.48}{\micro\watt}\\
    Delay & \SI{1016}{\pico\second} & \SI{563}{\pico\second} & \SI{1006}{\pico\second}\\
    PDP & \SI{5.93}{\femto\joule} & \SI{2.04}{\femto\joule} & \SI{4.51}{\femto\joule}\\
    Area & \SI{2.40}{\micro\meter\squared} & \SI{1.00}{\micro\meter\squared} & \SI{1.34}{\micro\meter\squared}\\
    Other & \SI{64.8}{\femto\coulomb} (\SI{116.7}{\femto\joule}) & \SI{40.3}{\femto\coulomb} (\SI{72.5}{\femto\joule}) & \SI{49.8}{\femto\coulomb} (\SI{89.6}{\femto\joule})\\
    \bottomrule
  \end{tabular}
\end{table}

\textbf{Supplementary comparison (conventional / BBM).} The $\Delta$ columns give the change of BBM relative to conventional (negative means BBM is lower).

\begin{table}[H]
  \centering\footnotesize
  \setlength{\tabcolsep}{3.2pt}
  \begin{tabular}{@{}p{3.2cm} c c c c r r r r@{}}
    \toprule
    Case & $T_p$ (ps) & $Q$ (fC/cycle) & $P_{avg}$ ($\mu$W) & PDP (fJ) & $\Delta Q$ & $\Delta$Power & $\Delta$Delay & PDP impr.\\
    \midrule
    Paper \SI{750}{\pico\second} sizing (delays not matched) & 507 / 563 & 67.3 / 40.3 & 6.06 / 3.63 & 3.07 / 2.04 & \SI{-40.2}{\percent} & \SI{-40.2}{\percent} & \SI{+11.2}{\percent} & \SI{33.5}{\percent}\\
    Re-optimised, target \SI{0.5}{\nano\second} & 495 / 505 & 44.1 / 38.3 & 3.97 / 3.45 & 1.96 / 1.74 & \SI{-13.1}{\percent} & \SI{-13.1}{\percent} & \SI{+2.0}{\percent} & \SI{11.4}{\percent}\\
    Re-optimised, target \SI{1.0}{\nano\second} & 1016 / 1006 & 64.8 / 49.8 & 5.84 / 4.48 & 5.93 / 4.51 & \SI{-23.2}{\percent} & \SI{-23.2}{\percent} & \SI{-0.9}{\percent} & \SI{23.9}{\percent}\\
    \bottomrule
  \end{tabular}
\end{table}

\textbf{Optimised parameter sets (nm):}
\begin{itemize}[nosep]
  \item 0.5\,ns conventional: \texttt{lp1=779, ln1=1113, wp2=249, lp2=632, wn2=371, ln2=1257}
  \item 0.5\,ns BBM: \texttt{lp1=913, wn1=239, ln1=196, wpt=260, lpt=189, lnf=283, lp2=197, ln2=593}
  \item 1.0\,ns conventional: \texttt{lp1=1330, ln1=1223, wp2=283, lp2=1168, wn2=784, ln2=1426}
  \item 1.0\,ns BBM: \texttt{lp1=1365, wn1=569, ln1=244, wpt=246, lpt=546, lnf=936, lp2=280, ln2=501}
\end{itemize}

\sect{9. OPTIMIZATION PERFORMED}
\textbf{Parameter Optimized:} Transistor sizing of both cells (channel lengths and widths; for BBM mainly $l_{P1}$, $w_{N1}$, $l_{N1}$, $w_{Pt}$, $l_{Pt}$, $l_{Nf}$, $l_{P2}$, $l_{N2}$ and the delay-element transistor $P_t$), for a fixed delay target, minimising the supply charge per cycle $Q$. Method: SciPy differential evolution driving ngspice (960 simulations per case).

\textbf{Initial Value:} Paper's \SI{750}{\pico\second} sizing: $Q=\SI{67.3}{\femto\coulomb}$ (conventional, $T_p=\SI{507}{\pico\second}$) and \SI{40.3}{\femto\coulomb} (BBM, $T_p=\SI{563}{\pico\second}$).

\textbf{Optimized Value:} $T_p\approx\SI{1.0}{\nano\second}$: $Q=\SI{64.8}{\femto\coulomb}$ (conventional) and \SI{49.8}{\femto\coulomb} (BBM). $T_p\approx\SI{0.5}{\nano\second}$: $Q=\SI{44.1}{\femto\coulomb}$ and \SI{38.3}{\femto\coulomb}.

\textbf{Reason:} The paper's own comparison sizes every cell for the same delay. Comparing the BBM cell at the paper's sizing with an un-resized conventional cell ($Q=\SI{67}{\femto\coulomb}$ vs \SI{40}{\femto\coulomb}, a misleading ``\SI{40}{\percent}'') is unfair, because the conventional chain at \SI{750}{\pico\second} sizing is far from the optimum for a \SI{0.5}{\nano\second} or \SI{1.0}{\nano\second} target (its charge falls to \SI{44}{\femto\coulomb} at \SI{0.5}{\nano\second}). Re-sizing both circuits for equal delay gives numbers I can stand behind. Limits: the optimiser is much smaller than the paper's (100-node SADE), so neither sizing is guaranteed to be globally optimal.

\sect{10. SIMULATION RESULTS}
Actual ngspice results (raw files in \texttt{results/}: \texttt{*.log} measurement logs, \texttt{*\_wave.csv} waveforms, \texttt{summary.json}, \texttt{summary\_opt.json}).

\begin{figure}[H]
  \centering
  \includegraphics[width=0.95\textwidth]{fig_stage3_currents.png}
  \caption{Stage-3 supply ($i_P$) and ground ($i_N$) currents. In the BBM cell the merged stage never carries a large overlapping current; the peak supply current of stage~3 is \SI{37.8}{\micro\ampere} (BBM) vs \SI{54.4}{\micro\ampere} (conventional).}
\end{figure}

\begin{figure}[H]
  \centering
  \includegraphics[width=0.72\textwidth]{fig_Q_vs_Tp_optimised.png}
  \caption{Charge per cycle at matched delay targets after re-optimisation (two targets only).}
\end{figure}

\begin{figure}[H]
  \centering
  \includegraphics[width=0.72\textwidth]{fig_Q_vs_Tp_scaling.png}
  \caption{One-parameter scaling sweeps from the paper's \SI{750}{\pico\second} sizing (not re-optimised; the conventional cell is \emph{not} at its optimum away from \SI{750}{\pico\second}, so this plot overstates the BBM advantage).}
\end{figure}

\textbf{Load sensitivity at the paper's \SI{750}{\pico\second} sizing} (delays not matched, so absolute savings are inflated; trend only).
\begin{table}[H]
  \centering\small
  \begin{tabular}{@{}r r r r@{}}
    \toprule
    $C_L$ & $T_p$ conv / BBM (ps) & $Q$ conv / BBM (fC) & BBM charge change\\
    \midrule
    \SI{0.5}{\femto\farad} & 497 / 554 & 64.8 / 38.0 & \SI{-41.4}{\percent}\\
    \SI{2}{\femto\farad}   & 507 / 563 & 67.3 / 40.3 & \SI{-40.2}{\percent}\\
    \SI{5}{\femto\farad}   & 522 / 578 & 72.6 / 45.3 & \SI{-37.6}{\percent}\\
    \SI{10}{\femto\farad}  & 543 / 597 & 81.6 / 54.0 & \SI{-33.9}{\percent}\\
    \bottomrule
  \end{tabular}
\end{table}

\textbf{Input-edge-rate sensitivity} (same sizing).
\begin{table}[H]
  \centering\small
  \begin{tabular}{@{}r r r@{}}
    \toprule
    Input edge & $T_p$ conv / BBM (ps) & $Q$ conv / BBM (fC)\\
    \midrule
    \SI{20}{\pico\second}  & 500 / 557 & 67.9 / 40.9\\
    \SI{50}{\pico\second}  & 507 / 563 & 67.3 / 40.3\\
    \SI{200}{\pico\second} & 533 / 584 & 67.9 / 40.7\\
    \SI{500}{\pico\second} & 562 / 606 & 70.5 / 43.2\\
    \bottomrule
  \end{tabular}
\end{table}

\textbf{Ablation} (documented, not used for the improvement figure). Removing the delay element (A and B joined by \SI{10}{\ohm}, $N_f$ deleted) with identical device sizes gives $T_p=\SI{155}{\pico\second}$ and $Q=\SI{28.9}{\femto\coulomb}$, a different (much faster) circuit, because the $P_t$ resistance is what creates the delay. It therefore cannot isolate the short-circuit saving. A simple $\min(i_P,i_N)$ crowbar estimate for stage~3 gives \SI{0.74}{\femto\coulomb} (conventional) and \SI{0.80}{\femto\coulomb} (BBM) per cycle, i.e.\ only a few percent of the stage-3 charge in this model.

\textbf{Required Results:}
\begin{itemize}[label={},leftmargin=1.2em,nosep]
  \item \ding{51}\ Functional waveform (Figs.~\ref{fig:inout}, \ref{fig:skew})
  \item \ding{51}\ Power measurement (Section~8 table)
  \item \ding{51}\ Delay measurement (Section~8 table)
  \item \ding{51}\ PDP calculation (Section~8 table)
  \item \ding{51}\ Optimization result (Section~9, Fig.~6)
  \item \ding{51}\ Comparative result (Section~8, Section~12)
\end{itemize}

\sect{11. RESULT}
The implemented circuits were successfully simulated (ngspice-42) and the performance was evaluated. The improvement is computed for the delay-matched re-optimised pair at $T_p\approx\SI{1.0}{\nano\second}$ (conventional \SI{1016}{\pico\second}, BBM \SI{1006}{\pico\second}).

\textbf{Power:} $P_{avg}=\SI{5.84}{\micro\watt}$ (conventional) vs \SI{4.48}{\micro\watt} (BBM) at \SI{50}{\mega\hertz}, \SI{23.2}{\percent} lower ($Q=\SI{64.8}{\femto\coulomb}$ vs \SI{49.8}{\femto\coulomb} per cycle).

\textbf{Delay:} \SI{1016}{\pico\second} (conventional) vs \SI{1006}{\pico\second} (BBM): \SI{0.9}{\percent} shorter, i.e.\ matched by design (target \SI{1.0}{\nano\second} $\pm\,\SI{2}{\percent}$).

\textbf{PDP:} \SI{5.93}{\femto\joule} (conventional) vs \SI{4.51}{\femto\joule} (BBM) ($\mathrm{PDP}=P_{avg}\times T_p$ at \SI{50}{\mega\hertz}).

\textbf{Transistor Count:} 8 (conventional) vs 10 (BBM).

\textbf{Improvement Achieved:} \SI{23.9}{\percent}\quad [$=(5.927-4.511)/5.927\times 100$, from the logged PDP values].

\sect{12. DISCUSSION}
At $T_p\approx\SI{1.0}{\nano\second}$ the BBM cell cuts PDP by \SI{23.9}{\percent} and charge by \SI{23.2}{\percent}. At about \SI{0.5}{\nano\second} the cuts are \SI{11.4}{\percent} (PDP) and \SI{13.1}{\percent} (charge). The paper's Table~1 shows the same pattern: a \SI{10}{\percent} saving at \SI{0.5}{\nano\second}, \SI{36}{\percent} at \SI{1.0}{\nano\second}, a bit over \SI{40}{\percent} from \SI{3.5}{\nano\second}, and no gain below about \SI{0.35}{\nano\second}. My saving at \SI{1}{\nano\second} is smaller than the paper's.

BBM is not free. It needs 10 transistors instead of 8, plus two wires between stages. At the paper's \SI{750}{\pico\second} sizing the BBM cell is also slower (\SI{563}{\pico\second} against \SI{507}{\pico\second}, with rise and fall delays of \SI{490}{\pico\second} and \SI{637}{\pico\second}), so the \SI{40}{\percent} charge saving at that point should not be quoted. The gate-area proxy at \SI{1.0}{\nano\second} is \SI{1.34}{\micro\meter\squared} for BBM and \SI{2.40}{\micro\meter\squared} for the conventional cell, but that is not a layout area. With sizes fixed, the gain shrinks as the load grows (\SI{41}{\percent} at \SI{0.5}{\femto\farad}, \SI{34}{\percent} at \SI{10}{\femto\farad}) and with slow input edges (\SI{38.7}{\percent} at \SI{500}{\pico\second}). I did not test delays at or below \SI{0.35}{\nano\second}, where the paper finds no advantage.

Why does it save charge? Eq.~(16) of the paper gives a charge proportional to $T_p$ with a smaller slope than Eq.~(9), because $Q_s$ is about zero. In my runs the crowbar charge of the conventional stage~3 is small, so I think most of the saving comes from how the delay is made. BBM gets it from a long, narrow $P_t$ (and long $P_1$ and $N_f$) that load only the small first stage. The conventional chain needs a big, long-channel inv2 whose capacitance is charged every cycle. The paper makes a similar remark about the third-stage gate capacitance. I did not separate the two effects rigorously.

\textbf{Limits.} Level-1 has no velocity saturation, sub-threshold current or gate leakage, so static power is zero, and only \SI{27}{\celsius} was run. $\lambda$, $\gamma$ and the capacitances are my assumptions, so absolute fC values cannot be compared with the paper's TSMC numbers. The unit-inverter sizes and $C_L$ are assumed too. There is one corner, no mismatch, a reduced optimiser, no layout parasitics (as in the paper), and the input driver's charge is not counted.

\sect{13. CONCLUSION}
I rebuilt the paper's break-before-make delay cell (PtNf) from its figures, wrote ngspice netlists for it and for a conventional 4-inverter reference, and simulated both with a substitute Level-1 \SI{0.18}{\micro\meter}-class model at \SI{1.8}{\volt}. Both work and give full-swing outputs. In the BBM cell the PMOS-gate and NMOS-gate nodes of the last stage switch \SIrange{155}{258}{\pico\second} apart (paper sizing), so the two transistors never conduct together.

With both cells re-sized for the same delay, BBM cut PDP by \SI{23.9}{\percent} at about \SI{1.0}{\nano\second} and \SI{11.4}{\percent} at about \SI{0.5}{\nano\second}, and the charge per cycle by \SI{23.2}{\percent} and \SI{13.1}{\percent}. That follows the paper's trend, more benefit at longer delay, but not its numbers. I would call it a qualitative confirmation for this substitute model, not a reproduction of the TSMC results.

\renewcommand{\refname}{14. REFERENCES}
\begin{thebibliography}{9}
  \bibitem{puhan} J.~Puhan, D.~Rai\v{c}, T.~Tuma, and \'A.~B\H{u}rmen, ``Break-before-Make CMOS Inverter for Power-Efficient Delay Implementation,'' \emph{The Scientific World Journal}, vol.~2014, Article ID 349131, 11~pp., 2014. doi:\href{https://doi.org/10.1155/2014/349131}{10.1155/2014/349131}. \url{https://pmc.ncbi.nlm.nih.gov/articles/PMC4265520/}
  \bibitem{ngspice} ngspice circuit simulator, version~42, \url{https://ngspice.sourceforge.io/} (simulations and \texttt{.meas} measurements).
  \bibitem{level1} Level-1 MOSFET model parameters taken from the caption of Fig.~4 of [1]. Unlisted parameters are my assumptions (see \texttt{models\_180nm\_L1.lib}).
\end{thebibliography}

\end{document}
